现在对算子 $G$ 应用 Fredholm 二择一定理 (Theorem~\ref{thm:fredholm-alternative-compact}), 得到以下结论.

\begin{Theorem}
\label{thm:elliptic-fredholm-alternative}
设 $L(\bm{x},D)$ 为 $2k$ 阶一致椭圆微分算子, 满足 Theorem~\ref{thm:weak-dirichlet-positive-shift-existence} 的假设. 那么对每个 $\mu\in\mathbb{C}$, Fredholm 二择一结论成立; 即以下两种情形之一成立:\begin{enumerate}\item 对每个 $f\in L^2(\Omega)$, 下列方程的 Dirichlet 问题都存在唯一弱解 $u\in H_0^k(\Omega)$:
\begin{equation}\label{eq:elliptic-fredholm-inhomogeneous-equation}
L(\bm{x},D)u-\mu u=f,
\end{equation}
即
\begin{equation}\label{eq:elliptic-fredholm-inhomogeneous-weak-equation}
B[v,u]-\mu(v,u)_{L^2(\Omega)}=(v,f)_{L^2(\Omega)}
\end{equation}
对所有 $v\in H_0^k(\Omega)$ 成立; 或者\item 至多存在有限个线性无关函数 $u_i\in H_0^k(\Omega)$, $i=1,\ldots,N$, 使得
\begin{equation}\label{eq:elliptic-fredholm-homogeneous-weak-equations}
B[v,u_i]-\mu(v,u_i)=0,
\end{equation}
对所有 $v\in H_0^k(\Omega)$ 成立.\end{enumerate}此外, 第二种情形成立时相应的值组成无穷离散集, 没有有限聚点.
\end{Theorem}

\begin{Proof}
首先将方程
\begin{equation}\label{eq:elliptic-fredholm-differential-eigenproblem}
Lu=\mu u
\end{equation}
写为
\begin{equation}\label{eq:elliptic-fredholm-shifted-eigenproblem}
(L+\widetilde\lambda)u=(\widetilde\lambda+\mu)u.
\end{equation}
然后作形式计算, 对 \eqref{eq:elliptic-fredholm-shifted-eigenproblem} 的两端作用 $G/\widetilde\lambda=(L+\widetilde\lambda)^{-1}$, 可见 \eqref{eq:elliptic-fredholm-differential-eigenproblem} 成立, 当且仅当 $u$ 解
\begin{equation}\label{eq:elliptic-fredholm-solution-operator-eigenrelation}
u=(L+\widetilde\lambda)^{-1}(L+\widetilde\lambda)u=\dfrac{\widetilde\lambda+\mu}{\widetilde\lambda}Gu=\dfrac{1}{\sigma}Gu.
\end{equation}
因此可见, $u\in L^2(\Omega)$ 为 $G$ 的对应于特征值 $\sigma$ 的特征函数, 当且仅当 $u$ 为 $L$ 的对应于特征值 $\mu$ 的特征函数, 其中
\begin{equation}\label{eq:elliptic-fredholm-sigma-from-mu}
\sigma=\dfrac{\widetilde\lambda}{\widetilde\lambda+\mu},
\end{equation}

\begin{equation}\label{eq:elliptic-fredholm-mu-from-sigma}
\mu=-\widetilde\lambda+\dfrac{\widetilde\lambda}{\sigma}.
\end{equation}
由 Fredholm 二择一定理, $G$ 的非零特征值具有有限重数, 因而 $L$ 的特征值也具有有限重数. 此外, $G$ 的特征值组成一个离散集, 其唯一可能的聚点为零; 又由于已使 $0$ 不是 $G$ 的特征值, 所以 $G$ 必有无穷多个特征值. 因而 $L$ 必有无穷多个特征值, 且没有有限聚点.

当 $\mu\ne-\widetilde\lambda$ 不是 $L$ 的特征值时, 注意到 $u\in H_0^k(\Omega)$ 为 \eqref{eq:elliptic-fredholm-inhomogeneous-equation} 的解, 当且仅当 $u$ 为下式的解:
\begin{equation}\label{eq:elliptic-fredholm-equivalent-compact-equation}
G(u)-\dfrac{\widetilde\lambda}{\mu+\widetilde\lambda}u=-\dfrac{1}{\mu+\widetilde\lambda}G(f).
\end{equation}
留待读者补充使这一形式论证成立所需的严密性. 唯一需要细致处理的是证明: 含 $G$ 的方程的解函数 $u$ (因而自然仅被视为 $L^2(\Omega)$ 中的函数), 实际上必须是嵌入 $L^2(\Omega)$ 中的 $H_0^k(\Omega)$ 中的函数 (从而可作为含 $L$ 的方程的弱解).
\end{Proof}
